\documentclass[10pt]{article}
\usepackage[T1]{fontenc}
\usepackage[utf8]{inputenc}
\usepackage[french]{babel}
\usepackage{amsmath,amssymb,mathrsfs,mathpazo}
\usepackage{geometry}
\geometry{a4paper,margin=2cm}
\usepackage{graphicx,array,tabularx,tabularray,colortbl,multirow,multicol,booktabs,enumitem,pifont,ulem,fancybox,diagbox,cancel}
\usepackage{pstricks,pstricks-add,pst-plot,pst-tree,pst-node,pst-text,pst-func,pst-math,pst-bspline,pst-3dplot,pst-eucl,pst-solides3d}
\let\pscalculate\undefined
\usepackage{pst-bezier}
\usepackage{tikz,pgfplots,tkz-tab}
\usetikzlibrary{arrows,arrows.meta,patterns,calc,positioning,trees,decorations.markings,shapes.geometric,shapes.misc,shapes.arrows,intersections,decorations.pathmorphing,decorations.pathreplacing,backgrounds,shadows,babel}
\pgfplotsset{compat=1.18}
\usepgfplotslibrary{fillbetween}
\usepackage[np]{numprint}
\usepackage{listings,makecell,pas-tableur,twoopt}
\let\barre\undefined
\usepackage[pstricks]{tablvar}
\definecolor{darkgreen}{RGB}{0,100,0}
\providecommand{\red}{\color{red}}
\providecommand{\blue}{\color{blue}}
\providecommand{\green}{\color{green}}
\providecommand{\black}{\color{black}}
\providecommand{\white}{\color{white}}
\providecommand{\gray}{\color{gray}}
\providecommand{\cyan}{\color{cyan}}
\providecommand{\magenta}{\color{magenta}}
\providecommand{\yellow}{\color{yellow}}
\providecommand{\violet}{\color{violet}}
\providecommand{\orange}{\color{orange}}
\providecommand{\euro}{\texteuro}
\providecommand{\EUR}{\texteuro}
\providecommand{\ieme}{\textsuperscript{e}}
\providecommand{\degres}{\ensuremath{^\circ}}
\providecommand{\degree}{\ensuremath{^\circ}}
\providecommand{\vv}[1]{\overrightarrow{#1}}
\providecommand{\Coord}[1]{#1}
\providecommand{\decosix}{}
\providecommand{\starredbullet}{\ensuremath{\bullet}}
\pagestyle{empty}
\setlength{\parindent}{0pt}
\usepackage{adjustbox}\DeclareUnicodeCharacter{2236}{:}\DeclareUnicodeCharacter{2190}{\ensuremath{\leftarrow}}
\providecommand{\R}{\mathbb{R}}
\providecommand{\N}{\mathbb{N}}
\providecommand{\Z}{\mathbb{Z}}
\providecommand{\Q}{\mathbb{Q}}
\providecommand{\D}{\mathbb{D}}
\providecommand{\C}{\mathbb{C}}
\providecommand{\e}{{\mathrm{e}}}
\providecommand{\vect}[1]{\overrightarrow{#1}}
\providecommand{\vectt}[1]{\overrightarrow{\mathrm{#1}}}
\providecommand{\vvt}[1]{\overrightarrow{\mathrm{#1}}}
\providecommand{\np}[2][]{{#2~#1}}
\providecommand{\vv}[1]{{\overrightarrow{#1}}}
\renewcommand{\d}{{\mathrm{d}}}
\providecommand{\degres}{{^\circ}}
\providecommand{\degre}{{^\circ}}
\providecommand{\blue}{\color{blue}}
\providecommand{\red}{\color{red}}
\providecommand{\enskip}{\quad}
\providecommand{\llbracket}{\lbrack\!\lbrack}
\providecommand{\rrbracket}{\rbrack\!\rbrack}
\providecommand{\euro}{{\text{€}}}
\providecommand{\no}{{\text{n°}}}
\providecommand{\interval}[2]{{[#1;#2]}}
\providecommand{\intervFO}[2]{{[#1;#2[}}
\providecommand{\dotfill}{\ldots}

\providecommand{\ds}{\displaystyle}
\providecommand{\textding}[1]{#1}
\providecommand{\Oij}{\(\mathrm{O};\vec{\imath},\vec{\jmath}\)}
\providecommand{\Oijk}{\(\mathrm{O};\vec{\imath},\vec{\jmath},\vec{k}\)}
\providecommand{\Ouv}{\(\mathrm{O};\vec{u},\vec{v}\)}
\providecommand{\intervFF}[2]{\ensuremath{[#1;#2]}}
\renewcommand{\euro}{\ifmmode\text{\texteuro}\else\texteuro\fi}
\DeclareUnicodeCharacter{00B0}{\ifmmode^{\circ}\else\textdegree\fi}
\DeclareUnicodeCharacter{20AC}{\ifmmode\text{\texteuro}\else\texteuro\fi}
% Énoncés : Éducation nationale. Transcription et corrigés : APMEP, Denis Vergès et contributeurs. Découpage et classement : CoopMaths, Sébastien Lozano, Éric Elter et contributeurs. Adaptation de la présentation : Maths Explorer.
% CC BY-SA : https://coopmaths.fr/www/about/
% https://coopmaths.fr/alea/static/e3c/2020/tex/e3c_2020_00_sujet13_4_cor.tex
\begin{document}
\noindent\textit{APMEP / CoopMaths. CC BY-SA. Présentation adaptée par Maths Explorer.}\par\medskip
\section*{Exercice 4 (5 points)}
	
	\subsection*{1. Étudier le signe de la fonction \(P\) définie sur \(\mathbb{R}\) par \(P(x) = x^2 + 4x + 3\).}
	
\[ P(x) = x^2 + 4x + 3 = (x + 2)^2 - 4 + 3 = (x + 2)^2 - 1 = (x + 3)(x + 1). \]
On sait que \(P(x) > 0\), sauf sur l’intervalle \(\left] -3, -1 \right[\) (-3 et -1 étant les racines du trinôme).

	\subsection*{2. Montrer que pour tout réel \(x\) de l’intervalle \(\left] -2, +\infty \right[\),
		\[ f'(x) = \dfrac{P(x)}{(x + 2)^2}, \]
		où \(f'\) est la fonction dérivée de \(f\).}
	
	Pour \(x > -2\), \(f\) est dérivable et sur \(\left] -2, +\infty \right[\) :
	\[ f'(x) = \dfrac{(2x + 1)(x + 2) - (x^2 + x - 1)}{(x + 2)^2} = \dfrac{2x^2 + 4x + x + 2 - x^2 - x + 1}{(x + 2)^2} = \dfrac{x^2 + 4x + 3}{(x + 2)^2} = \dfrac{P(x)}{(x + 2)^2}. \]
	
	\subsection*{3. Étudier le signe de \(f'(x)\) sur \(\left] -2, +\infty \right[\) et construire le tableau de variations de la fonction \(f\) sur \(\left] -2, +\infty \right[\).}
Comme \((x + 2)^2 > 0\) sur \(\left] -2, +\infty \right[\), le signe de \(f'(x)\) est celui du numérateur \(P(x)\) étudié à la question 1.
Donc sur \(\left] -2, -1 \right[\), \(f'(x) < 0\) : la fonction est décroissante et sur \(\left] -1, +\infty \right[\) la fonction \(f\) est croissante.
	
	\begin{center}
	\begin{tikzpicture}[double distance=2pt]
		\tkzTabInit{$x$/1,$f'(x)$/1,$f(x)$/2}{$-\infty$,$-1$, $+\infty$}
		\tkzTabLine{,-,z, +}
		\tkzTabVar{+/,-/$-1$,+/}
	\end{tikzpicture}
\end{center} 
	
	\subsection*{4. Donner le minimum de la fonction \(f\) sur \(\left] -2, +\infty \right[\) et la valeur pour laquelle il est atteint (on donnera les valeurs exactes).}
	
	D’après la question précédente \(f(-1) = \dfrac{1 - 1 - 1}{-1 + 2} = -1\) est le minimum de la fonction \(f\) sur \(\left] -2, +\infty \right[\).
	
	\subsection*{5. Déterminer le coefficient directeur de la tangente \(T\) à la courbe \(C_f\) au point d’abscisse 2.}
	
	Le coefficient directeur de la tangente \(T\) à la courbe \(C_f\) au point d’abscisse 2 est le nombre dérivé \(f'(2)\) :
	\[ f'(2) = \dfrac{4 + 8 + 3}{(2 + 2)^2} = \dfrac{15}{16} = 0,9375. \]

\end{document}
